\section{Discussion}
\label{sec:discussion}

\subsection{Why the methods behave as observed}
The fourth-order formulation creates an explicit property mode and spatial--property measurement operator. Matrix unfolding is lossless; the distinction from independent 3D is the imposed sharing, not avoidance of information loss during reshaping. E8 correctly found that rigid joint coefficients lose to independent property models. E9 shows why that result cannot be generalised to every fourth-order architecture: jointly estimated spatial factors with separate property coefficients improve many matched sparse comparisons, whereas the rigid model still fails. Property rank two adds no property compression; the gain comes from reusing spatial factors and regularising their estimation. Property rank one discards too much property-specific variation and is never selected for the rigid model.

The ensemble prior removes the common response, leaving a field-wide pressure shift and near-well deviations (Supplementary Section~S4). Energy weighting shrinks weak POD directions in underdetermined recovery; spatial truncation instead imposes an irreducible representation floor. Existing wells may leave the sampled basis nearly singular, so regularization trades variance for bias. Factor storage savings do not imply lower peak process memory because this implementation materialises the basis matrix.

\subsection{Implications for reservoir monitoring}
The benchmark gives four practical lessons. First, compare any sensing method with a no-measurement reference: persistence is strong in quasi-steady windows, and the P2 prior is exact initially. Second, a simulation-ensemble basis and existing-well measurements can recover the scenario-specific pressure deviation, but prior-plus-interpolation remains a necessary reference. Third, well ranking mainly protects unregularised least squares; it matters less with $\ell_1$ estimation. Fourth, saturation channels add little for this fixed-geology ensemble, and pressure-only fitting degrades saturation reconstruction. Cross-property inference therefore requires separate validation.

Grid-wide channels are idealised cell observations, not physical sensors or a performance bound for new wells. Their concentration near wells and high-variability regions (Supplementary Section~S9) only describes this ensemble.

\subsection{What optimization changes}
The nested study separates four conclusions that the fixed-rank benchmark could not. First, spatial rank is the dominant original TBMD limitation: solver and placement optimization on the old basis barely changes its all-well error, whereas rank expansion removes most of the representation floor. Second, completely independent property factors can give the best full-state equal-property objective, but create a poorly conditioned block basis under sparse sensing. Third, E9's shared-factor/independent-head model supplies a middle ground: it preserves property-specific coefficients while pooling spatial-factor estimation, with its clearest gain at five wells and 30 common grid channels. Fourth, tensor compression is real but not free. The compressed independent basis stores \OptCompressionCompression{} times fewer values than POD-E and has the better joint objective at 300 grid channels, while POD-E still has substantially lower pressure error and remains the preferred method for existing joint wells. These are architecture- and regime-specific statements, not overall leaderboard wins.

\subsection{Generality and relation to previous work}
The construction applies to co-registered variables: properties enlarge mode 3 without changing the derivation. Proposition~\ref{prop:tbmd_pod} shows why tensor order alone cannot guarantee an accuracy gain; ranks, sharing and atom scaling determine the representation. E9's pressure--saturation anomalies contain one strongly aligned global direction and many weakly aligned spatial directions, explaining why partial sharing succeeds where common coefficients do not. Gains reported elsewhere \cite{zhong2024tbmd,farazmand2023} may occur for different data structure or truncation. The Brugge conclusions require validation on other properties, grids and geological ensembles.

\section{Limitations}
\label{sec:limitations}
The study uses one reservoir model, an areal two-dimensional export, and ten control scenarios with fixed geology; it excludes geological uncertainty, three-dimensional layering and longer operating histories, and the leave-one-out folds are not independent reservoir samples. Time is available only as snapshot indices, and the P2 prior assumes time-aligned training-ensemble trajectories. Measurements are idealised cell values: well-bore effects, gauge drift, correlated errors and saturation-measurement difficulty are not modelled beyond additive Gaussian noise; E9's headline comparison is noise free. The optimization is nested within P2, but the Brugge benchmark and prior outer results had already informed the diagnosis and bounded search; the outer folds are untouched within this run, not an externally preregistered data set. Anomaly-relative error and mask-aware SSIM are declared scientific choices; another application may prefer a different reference state, window or property weighting. The exploratory coupling--gain associations have only ten paired folds and do not identify a causal threshold. Model storage is measured exactly, whereas method-specific peak process memory is not. Kriging, neural decoders \cite{fukami2021,zou2025prosnet} and ensemble data assimilation \cite{evensen2009} were not included as baselines. Unrestricted public redistribution of the transformed TNO-based state exports has not been established; the corresponding author can instead provide them upon reasonable request, subject to applicable data-use and access conditions. Released result files support checking every table, statistic and quoted number (data availability statement).
